% Generated by experiments/scripts/export_latex.py -- do not edit by hand.
% Source: experiments/data/2-processed-tables-and-evaluations/revised-usd/window-w1/eval_summary.json
\begin{table}[htbp]
\centering
\caption[허락을 받는 트레이더와 라벨이 있는 다른 모든 트레이더의 비교.]{동결일에 받은 허락 수가 0보다 큰 GMX~V2 계약 트레이더(받는 쪽)를, 라벨이 있는 다른 모든 트레이더와 두 기간에서 비교했어요(W1의 트레이더 $n=70$, W0의 트레이더 $n=88$). 받는 쪽과 다른 트레이더의 비교, 그리고 영향 점검은 표~\ref{tab:neutral-label-contrasts}의 대비와 함께 등록했어요. 여기서는 고친 점수에 적용했어요. 트레이더 가운데 가장 큰 받은 허락 수는 W1에서 1, W0에서 3이에요. 그래서 영향 점검에서 빼는 받는 쪽 열 개는 같은 값(동점) 가운데서 점(노드)의 순서대로 골랐어요. 영향 점검은 남은 트레이더에 대해 계약을 단위로 한 짝지은 부트스트랩이에요(다시 뽑기 2{,}000번, 시드 42).}
\label{tab:neutral-label-recipients}
\fitwidth{%
\begin{tabular}{lcc}
\toprule
 & 받는 쪽 & 다른 트레이더 \\
\midrule
\multicolumn{3}{l}{\emph{기간 W1, 2026년 7월--9월 (트레이더 $n=70$)}} \\
트레이더 수 & 37 & 33 \\
청산이 없는 비율 & 1.000 & 0.727 \\
청산 없이 닫은 비율의 평균 & 1.000 & 0.878 \\
이익을 낸 비율의 평균 & 0.325 & 0.479 \\
손실이 아닌 비율의 평균 & 0.433 & 0.644 \\
라벨 기간에 닫은 횟수의 중앙값 & 24.0 & 11.0 \\
\multicolumn{3}{l}{받은 허락 수가 가장 큰 받는 쪽 열 개를 뺀 P: $\Delta\tau=0.073$ [-0.039, 0.186]} \\
\midrule
\multicolumn{3}{l}{\emph{기간 W0, 2026년 4월--6월 (트레이더 $n=88$)}} \\
트레이더 수 & 40 & 48 \\
청산이 없는 비율 & 1.000 & 0.750 \\
청산 없이 닫은 비율의 평균 & 1.000 & 0.919 \\
이익을 낸 비율의 평균 & 0.362 & 0.474 \\
손실이 아닌 비율의 평균 & 0.402 & 0.607 \\
라벨 기간에 닫은 횟수의 중앙값 & 31.0 & 10.0 \\
\multicolumn{3}{l}{받은 허락 수가 가장 큰 받는 쪽 열 개를 뺀 P: $\Delta\tau=0.188$ [0.046, 0.331]} \\
\bottomrule
\end{tabular}
}
\end{table}
